\scanpage{6}
\mspage{(4)}
Donc on aurait envie de prendre
\[
\deleted{\mathbb Q}\ \mathbb Z(\{X_i\}[dX_i])_{0\leq i\leq n}/(\textstyle\sum X_i-1,\sum dX_i)_{\mathrm{pd}}
\]
où $\{\ \}$ désigne les pol. à p.d., mais ce n'est plus
\[
\simeq\mathbb Z\{X_0,\ldots,X_{n-1}\}[dX_0,\ldots,dX_{n-1}]
\]
(\deleted{qui est} \inserted{bien} \deleted{qui est} une résolution de $\mathbb Z$), car si on veut tirer de la relation $\sum_0^n X_i-1=0$,
\[
X_n=1-\sum_0^{n-1}X_i\quad
(\inserted{\text{et}}\ dX_n=-\sum_0^{n-1}dX_i\ \text{de }\sum_{\unsure{1}}^n dX_i=0)
\]
-- on a le ``\unsure{hic|bec}'' que $1-\sum_0^{n-1}X_i$ n'appartient pas à l'idéal à p.d. donné dans $\mathbb Z\{X_0,\ldots,X_{n-1}\}$ (donc on ne sait pas comment envoyer $\mathbb Z\{X_0,\ldots,X_{\unsure{n-1|n}}\}$ dans $\mathbb Z\{X_0,\ldots,X_{n-1}\}$ avec \unsure{même} \deleted{\unsure{noyau}} $\sum_0^n X_i-1$ dans le noyau....).

On s'en tire en prenant un anneau de coeff. différent de $\mathbb Z$, soit $S$, avec une \deleted{indéterminée} \inserted{``paramètre''} $t\in S$ fixé dont on \uline{sache prendre des p.d.} \inserted{(i.e. $t\in J$, $J$ idéal à p.d. de $S$)}, et en remplaçant les équations $\sum X_i=1$ de $E^{[n]}$ par l'équation
\[
\sum X_i=t\qquad\text{dans }S^{n+1}
\]
et définissant
\[
\unsure{(C^\bullet_{\mathrm{DRpd}\,t_{[n]}}(S,t))}=\left(S\{X_i\}_{0\leq i\leq n}[dX_i]_{0\leq i\leq n}\right)/(\textstyle\sum X_i-t,\sum dX_i)_{\mathrm{pd}}
\]
(On divise par l'\uline{idéal à p.d.} engendré par $\sum X_i-t$, $\sum dX_i$, ce qui a un sens
\mspage{5}
